%\textit{To conclude, I will present which method performed better and the consequences of having more stable portfolios: less rebalancing necessary, lower transaction costs, etc. I would also like to shortly discuss what the alternative methods can still not achieve and some ideas to further improve López de Prado's method. However, I will keep this very last part and won't get into details of how to solve the problems or how to precisely improve López de Prado's algorithm.}

Asset managers confronted with the asset allocation problem in section \ref{section:TheAssetAllocationProblem} can choose among a huge number of asset allocation models. In this thesis, we have described optimal portfolios computed using Markowitz's portfolio allocation model, and focused on analysing its risk-based version, i.e. Markowitz's minimum volatility portfolios. Markowitz's minimum volatility portfolios are the portfolios that showed the lowest variance during the time window from which past data was sampled to compute it (past $T$ periods returns used to compute the sample covariance matrix of asset returns). These portfolios are however highly sensitive to the condition number of the sample covariance matrix, which rises when correlations rise and the differences in return's variances of the assets become smaller. We have shown that this can lead to a significant under-performance in terms of out-of-sample risk-adjusted returns, concentration and turn-over of Markowitz's minimum volatility portfolios compared to naive portfolio allocation models like the equal weight and inverse variance allocation models.

To deal with instability in Markowitz's minimum volatility portfolios, we have described two methods that act on the covariance matrix used as input in the optimisation problem and one method which is a hybrid between the inverse variance allocation and Markowitz's minimum volatility allocation. The first two are two covariance shrinkage methods developed by Ledoit and Wolf and another developed by De Nard, while the last is a method developed by López de Prado. We have also tested these 3 methods in the one-year period of market stress given by the outbreak of the COVID-19 crisis in 2020, and compared them to the original Markowitz minimum volatility model and the two naive allocation models mentioned above (equal weight and inverse variance allocation models).

Covariance shrinkage methods try to tackle instability by defining a structured predictor for the covariance matrix, and using a convex combination of the structured predictor and the sample covariance matrix as input for the computation of Markowitz's minimum volatility portfolios. The different covariance shrinkage methods are given by different definitions of the structured predictor and of the shrinkage constant that defines the factor with which the structured predictor is weighted in the convex combination.

Ledoit and Wolf propose using the constant correlation matrix as structured predictor and the shrinkage constant given by the consistent estimator of the shrinkage constant that minimises the loss function given by the Frobenius norm of the difference between the shrunk and true covariance matrix and trimmed to the interval $[0,1]$. The resulting shrunk covariance matrices exhibited a lower condition number in our practical implementation. Moreover, using the resulting shrunk covariance matrix for the computation of Markowitz's minimum volatility portfolios also led to sharply improved out-of-sample risk-adjusted returns while reducing portfolio turnover. However, this covariance shrinkage method failed to reduce portfolio concentration in our practical implementation. Ledoit and Wolf's method comes with some limitations, which are given by the fact that all assets are considered to be equal, and no differentiation between asset classes, sectors or regions are made in the computation of the structured predictor. Moreover, the consistent estimator of the optimal shrinkage constant is only consistent if we assume that asset return's random variables are independent and identically distributed with finite fourth moments, which is not true in general.

De Nard's covariance shrinkage method tries to tackle the limitations of Ledoit and Wolf's method by clustering the assets into homogeneous groups to later define a structured predictor of the covariance matrix for each group separately. To do this, the K-means algorithm is used, and a new structured predictor called the generalised constant-variance-covariance shrinkage target is defined. The shrinkage constant in De Nard's method is the same as in Ledoit and Wolf's method, but using the generalised constant-variance-covariance shrinkage target as input for its computation. In our practical implementation, this method led to the lowest input covariance matrix' condition numbers, which were well below those of the sample covariance matrices and those of Ledoit and Wolf's shrunk covariance matrices. The condition numbers were also very stable across time, and a significant reduction in concentration could be observed. However, this method resulted in the highest turnover figures and was the worst performing in terms of out-of-sample risk-adjusted returns among all methods simulated. This might be explained by the fact that our practical implementation only used 20 assets, which were clustered into 4 homogeneous groups. Such a reduced number of assets and the fact that the K-means algorithm is very sensitive to outliers, could be the reason for the bad results in terms of Sharpe ration and portfolio turnover.

Finally, we described and implemented López de Prado hierarchical risk parity model, which can be considered a hybrid method between the inverse variance allocation and Markowitz's minimum volatility models. It was developed with the objective of improving the stability and out-of-sample performance of Markowitz's minimum volatility portfolios. It, first, hierarchically classifies the assets in the allocation problem using the sample covariance and its corresponding correlation matrix to give the problem a structure. Then, it allocates capital along the hierarchy using the inverse variance allocation model between the clusters. This method is not only an asset allocation model, but can also be understood as a framework that allows for some flexibility in terms of defining the tree structure and how allocations are then split between clusters in the structure. In our practical implementation, López de Prado hierarchical risk parity method preformed the best with respect to out-of-sample risk-adjusted returns. With respect to portfolio turnover, this method was only second to the equal weight allocation method. However, hierarchical risk parity portfolios were consistently among the most concentrated portfolios across time.

We conclude that asset managers should not discard naive portfolio allocation models like the equal weight or inverse variance allocation models in favour of more complex methods like Markowitz's. If they are to consider using Markowitz's minimum volatility portfolios, they must keep in mind the sensitivity to the sample covariance matrix' condition number in their computation. Covariance shrinkage methods can indeed reduce the condition number of the input covariance matrix and lead to better performance in terms of of risk-adjusted returns, turnover and concentration. Finally, alternative asset allocation methods like the hierarchical risk parity method can be good substitutes for Markowitz's model, even if the portfolios generated are not optimal in-sample.

As a final note, we would like to remind again that the empirical results of the practical implementation are limited by the fact that we have only analysed one past period of market stress. As it is often quoted in financial literature and marketing materials, past returns are not a guarantee of future performance. Moreover, better performance and optimal portfolios in periods of market stress might not result in better performance and optimal portfolios outside periods of market stress, and the overall (periods composed of market stress and calm markets phases) performance, as well as optimal portfolios, might differ from those in market stress phases.